Large-scale aerospace simulation environments generate substantial volumes of raw data
whose post-processing represents a critical bottleneck in analytical workflows. When
simulation batches are processed sequentially on a single machine, the time required to
derive meaningful metrics and aggregate values grows proportionally to the number of
executions, limiting the analytical agility required for Capability-Based Planning in
defense contexts.

This work proposes, implements, and evaluates a distributed computing infrastructure designed
to accelerate the post-processing stage of scenario-based aerospace simulations produced by
the ASA project, developed at the Institute for Advanced Studies of the Brazilian Air Force.
The project, titled AeroSOuDA, combines Apache Spark (via PySpark) for parallel metric
extraction with the Hadoop Distributed File System (HDFS) for scalable storage with data
locality, exposing an interactive command-line interface through which analysts define metric
and aggregation functions before submitting the job, without requiring advanced computational
knowledge.

Two deployment modes are presented: a local cluster via Docker Compose for prototyping,
and a remote virtual machine deployment accessed by a Python client over SSH and SFTP. The
experimental evaluation in two analytical scenarios over two independent datasets (1,3\,GB
and 1,5\,GB, respectively) demonstrates speedups of $2{,}86\times$ and $3{,}14\times$ with
four executors in a single-host cluster, corresponding to parallel efficiencies of
$71{,}4\%$ and $78{,}6\%$. A scalability experiment with one, two, and four executors on
the larger dataset confirms that the architecture scales beneficially, reaching $93{,}4\%$
of efficiency with two executors. The SSH communication overhead remains below $1\%$ of the
total execution time in both campaigns, confirming the feasibility of the remote deployment
model.

The results establish that the MapReduce paradigm, implemented via PySpark and HDFS,
provides a practical and scalable foundation for the distributed post-processing of
simulation batches in an aerospace defense context, reducing execution time by $65$--$68\%$
relative to the sequential baseline while keeping the raw dataset entirely on the server.
The measured speedup increases with the size of the dataset, consistent with the predictions
of Amdahl's Law derived from the scalability experiment.